\section{Testing and Validation}
\label{sec:testing}
%% ===================================================================

\FractalShark{} maintains three independent test projects, each
targeting a different layer of the system.  Together they provide
coverage from low-level GPU arithmetic through high-level rendering
correctness.

\subsection{CrummyTest (rendering pipeline)}
\label{subsec:crummytest}

\code{CrummyTest} (in \code{FractalSharkLib/CrummyTest.cpp}) is a
functional test suite invoked from the application's right-click menu
(\emph{Run Basic Test}).  It exercises the full rendering pipeline by
calling \code{Drain()} followed by \code{CalcFractal(true)}---the
direct rendering path (\cref{subsec:two-render-paths})---so that
iteration results are written to \code{m\_CurIters} and can be
inspected or saved.

The suite iterates over every \code{RenderAlgorithm} entry, rendering
multiple built-in views (including Views~0, 5, 10, and~11) with both
32-bit and 64-bit iteration types.  Eight named tests run in sequence:

\begin{itemize}
  \item \code{TestBasic}---renders each algorithm at several views and
        iteration types.
  \item \code{TestReferenceSave}---saves and reloads reference orbits,
        verifying round-trip correctness.
  \item \code{TestVariedCompression}---tests reference-orbit compression
        at varying error thresholds.
  \item \code{TestImaginaLoad}---loads orbits stored in Imagina format.
  \item \code{TestPerturbedPerturb}---exercises perturbed iteration
        chains for numerical correctness.
  \item \code{TestWindowResize}---resizes the application window and
        verifies re-rendering completes correctly.
  \item \code{TestStringConversion}---validates string-to-precision and
        precision-to-string conversion utilities.
  \item \code{TestGrowableVector}---validates the
        \code{GrowableVector} data structure
        (\cref{sec:disk_backed_growable_vectors}).
\end{itemize}

\noindent
A separate stress test, \code{TestReallyHardView27}, renders View~\#27
(period ${\sim}28$~billion, one quadrillion iterations) and is intended
as an overnight correctness and stability check rather than a routine
regression test.

\subsection{HpSharkFloatTest (GPU arithmetic)}
\label{subsec:hpsharkfloattest}

\code{HpSharkFloatTest} is a standalone CUDA executable (in
\code{HpSharkFloatTest/}) that validates the GPU high-precision
reference-orbit implementation described in \cref{subsec:ref-orbit-gpu}.
Its sole operator is \code{ReferenceOrbit2}, which exercises the complete
fused v2 recurrence.

Each vector is validated through complementary comparisons:

\begin{enumerate}
  \item \textbf{GPU vs.\ MPIR.}  The fused GPU result is converted to
        MPIR values and compared with the same recurrence evaluated in
        arbitrary-precision arithmetic.  MPIR is the ground truth.
  \item \textbf{CPU oracle vs.\ MPIR.}  The CPU
        \code{ReferenceReferenceOrbit2} implementation uses the
        \code{HpSharkFloat} representation and is independently checked
        against MPIR.
  \item \textbf{Pipeline checksums.}  Deterministic host and device
        checksums compare corresponding fused-stage boundaries and identify
        the first stage at which their intermediate representations diverge.
\end{enumerate}

\noindent
When a test fails, the checksum-guided debugging infrastructure
(\cref{sec:checksums}) reduces the mismatch to a targeted fused pipeline
stage.  Hardcoded and generated vectors cover signs, zero values,
carry/borrow behavior, alignment shifts, and fused-stage bounds.  The
Newton--Raphson cases use the same Ref2 engine to exercise first-derivative
and checkpoint-resume paths.  Performance benchmarks exercise the full
reference-orbit pipeline at three zoom depths: View~5
(${\sim}16{,}000$ iterations), View~30 (${\sim}670{,}000$ iterations), and
View~32 (${\sim}30{,}000{,}000$ iterations).

\subsection{FractalSharkTest (utility types)}
\label{subsec:fractalsharktest}

\code{FractalSharkTest} (in \code{FractalSharkTest/}) is a standalone
CPU executable that validates the fundamental data types and conversion
utilities used throughout the codebase.  It uses a lightweight
assertion-based framework (\code{TEST}, \code{ASSERT\_EQ},
\code{ASSERT\_NEAR}) and covers:

\begin{itemize}
  \item \code{HDRFloat<T>} arithmetic and normalization
        (\cref{sec:hdr-float}).
  \item \code{HDRFloatComplex} operations (complex multiply, add,
        magnitude).
  \item \code{HighPrecision} string parsing and round-trip conversion.
  \item \code{PointZoomBBConverter} coordinate transforms
        (\cref{subsec:coordinates}).
\end{itemize}

The tiered allocation strategy introduced earlier in this chapter addresses one aspect of memory
pressure.  The next chapter examines a complementary mechanism---disk-backed
growable vectors---that manages the large reference-orbit buffers whose size
can exceed available physical memory.
